% ============================================================
% preamble.tex — Cấu hình toàn cục cho tiểu luận PSO–FCM đánh giá rủi ro tín dụng
% Engine: XeLaTeX (BẮT BUỘC — KHÔNG dùng pdfLaTeX)
% ============================================================

% ---- Font & Tiếng Việt ----
\usepackage{fontspec}
% Font nhúng trong thư mục fonts/ của project (không phụ thuộc font hệ thống).
% Nhờ vậy compile được ở mọi môi trường: máy local, Docker TeX Live, Overleaf.
% Đường dẫn Path tính từ thư mục gốc project (nơi chạy lệnh compile main.tex).
\setmainfont{Times New Roman}[
  Path = fonts/,
  Extension = .ttf,
  UprightFont = times,
  BoldFont = timesbd,
  ItalicFont = timesi,
  BoldItalicFont = timesbi
]
% Sans-serif không dùng trong tài liệu → không khai báo Arial (đỡ 1 font phải
% nạp + dựng cache, giúp Overleaf compile lần đầu nhanh hơn, tránh timeout).
\setmonofont{Courier New}[
  Path = fonts/,
  Extension = .ttf,
  Scale = 0.9,
  UprightFont = cour,
  BoldFont = courbd,
  ItalicFont = couri,
  BoldItalicFont = courbi
]

\usepackage{polyglossia}
\setdefaultlanguage{vietnamese}
% polyglossia tự load hyphenation patterns tiếng Việt

% ---- Cỡ chữ 13pt (chuẩn báo cáo khoa học VN) ----
% Lớp report chỉ hỗ trợ tối đa 12pt; XeLaTeX dùng font TrueType scalable nên
% có thể đặt cỡ tùy ý. Đặt thân bài về 13pt, baselineskip 15.6pt (sẽ được
% setspace nhân 1.5 ở dưới).
\AtBeginDocument{\fontsize{13pt}{15.6pt}\selectfont}

% ---- Kích thước trang & lề ----
\usepackage[
  a4paper,
  left=3cm,
  right=2cm,
  top=2.5cm,
  bottom=2.5cm
]{geometry}

% ---- Giãn dòng ----
\usepackage{setspace}
\onehalfspacing   % giãn dòng 1.5

% ---- Tiêu đề, mục lục, đánh số ----
\usepackage{titlesec}
% Định dạng chapter: căn giữa, in hoa, 14pt đậm
\titleformat{\chapter}[display]
  {\normalfont\fontsize{14pt}{18pt}\selectfont\bfseries\centering}
  {\MakeUppercase{\chaptertitlename}\ \thechapter}
  {0pt}
  {\MakeUppercase}
\titlespacing*{\chapter}{0pt}{-10pt}{20pt}
% Mục và tiểu mục cùng cỡ thân bài (13pt), in đậm; tiểu mục thêm nghiêng.
\titleformat{\section}
  {\normalfont\fontsize{13pt}{15.6pt}\selectfont\bfseries}{\thesection}{0.6em}{}
\titleformat{\subsection}
  {\normalfont\fontsize{13pt}{15.6pt}\selectfont\bfseries\itshape}{\thesubsection}{0.6em}{}

% ---- Dấu chấm dẫn (dot leader) cho mục CHƯƠNG trong mục lục ----
% Mặc định mục cấp chương không có dấu chấm dẫn → tiêu đề ngắn để lại khoảng
% trống lớn tới số trang. titletoc (đi kèm titlesec) thêm dấu chấm cho đồng
% nhất với mục con.
\usepackage{titletoc}
% Mục lục in "CHƯƠNG 1." trước tên chương; main.tex đổi tiền tố thành
% "PHỤ LỤC" sau \appendix.
\newcommand{\chaptertocprefix}{CHƯƠNG}
\titlecontents{chapter}[0pt]
  {\addvspace{6pt}\bfseries}
  {\chaptertocprefix\ \thecontentslabel.\quad}
  {}
  {\titlerule*[0.6pc]{.}\contentspage}

% ---- Toán học ----
\usepackage{amsmath}
\usepackage{mathtools}
% Font toán cùng họ Times với thân bài (TeX Gyre Termes Math có sẵn trong
% TeX Live và Overleaf). unicode-math đã gồm các ký hiệu của amssymb.
\usepackage{unicode-math}
\setmathfont{texgyretermes-math.otf}

% ---- Hình ảnh ----
\usepackage{graphicx}
\graphicspath{{hinh/}}   % ảnh đặt trong thư mục hinh/

% ---- Vẽ khung trang bìa ----
\usepackage{tikz}        % khung trang trí trang bìa (remember picture/overlay)

% ---- Bảng biểu ----
\usepackage{booktabs}
\usepackage{array}
\usepackage{longtable}   % bảng dài nhiều trang

% ---- Caption cho hình và bảng ----
\usepackage{caption}
\captionsetup{
  font=small,
  labelfont=bf,
  labelsep=period
}

% ---- Mã giả thuật toán ----
% Dùng cặp package algorithm + algorithmic, có sẵn trong mọi bản TeX Live nên
% biên dịch được cả trên Docker lẫn Overleaf mà không cần cài thêm.
\usepackage{algorithm}
\usepackage[noend]{algorithmic}  % noend: khối lồng nhau thể hiện bằng thụt lề, không in dòng "end"
\floatname{algorithm}{Thuật toán}
\renewcommand{\algorithmicrequire}{\textbf{Đầu vào:}}
\renewcommand{\algorithmicensure}{\textbf{Đầu ra:}}
% Từ khoá mã giả tiếng Việt
\renewcommand{\algorithmicfor}{\textbf{với}}
\renewcommand{\algorithmicforall}{\textbf{với mọi}}
\renewcommand{\algorithmicdo}{\textbf{thực hiện}}
\renewcommand{\algorithmicwhile}{\textbf{trong khi}}
\renewcommand{\algorithmicif}{\textbf{nếu}}
\renewcommand{\algorithmicthen}{\textbf{thì}}
\renewcommand{\algorithmicelse}{\textbf{ngược lại}}
\renewcommand{\algorithmicrepeat}{\textbf{lặp}}
\renewcommand{\algorithmicuntil}{\textbf{cho tới khi}}
\renewcommand{\algorithmicreturn}{\textbf{trả về}}
\renewcommand{\algorithmicto}{\textbf{tới}}
\renewcommand{\algorithmicand}{\textbf{và}}
\renewcommand{\algorithmicor}{\textbf{hoặc}}

% ---- Màu sắc & liên kết ----
\usepackage{xcolor}
\usepackage[
  colorlinks=true,
  linkcolor=black,
  citecolor=black,
  urlcolor=blue,
  bookmarks=true,
  bookmarksnumbered=true,
  pdfencoding=auto
]{hyperref}

% ---- Tài liệu tham khảo (BibLaTeX + Biber) ----
% style=numeric-comp: đánh số [1] và gộp dải số liên tiếp, gần với chuẩn IEEE.
\usepackage[
  backend=biber,
  style=numeric-comp,
  sorting=none,
  maxbibnames=99,
  giveninits=true
]{biblatex}
% polyglossia đặt ngôn ngữ chính là vietnamese, nhưng biblatex không có
% vietnamese.lbx → các từ nối ("in", "and", "pages"...) bị lỗi dính chữ
% ("inIEEE Transactions..."). Ánh xạ sang bibstring tiếng Anh để khắc phục.
\DeclareLanguageMapping{vietnamese}{english}
\addbibresource{noi-dung/tai-lieu-tham-khao.bib}
% Trường note trong .bib chứa ghi chú nội bộ của nhóm (số bài, đường dẫn PDF,
% lý do thay bài) → bỏ khi in danh mục tài liệu tham khảo.
% URL dạng https://doi.org/... trùng với DOI → bỏ URL khi mục đã có DOI.
\DeclareSourcemap{
  \maps[datatype=bibtex]{
    \map{\step[fieldset=note, null]}
    \map{
      \step[fieldsource=doi, final]
      \step[fieldset=url, null]
    }
  }
}

% ---- Nhúng mã nguồn (phụ lục) ----
% Dùng fancyvrb (\VerbatimInput) thay cho listings: XeLaTeX xử lý UTF-8 native
% nên giữ ĐÚNG dấu tiếng Việt trong comment code. (listings băm theo byte →
% hỏng dấu tổ hợp tiếng Việt.)
\usepackage{fancyvrb}
% fancyvrb không tự bẻ dòng dài → dùng cỡ chữ nhỏ (\scriptsize) để dòng code
% dài nhất (~96 ký tự) vừa trong khung.
\fvset{
  frame=single,
  numbers=left,
  numbersep=5pt,
  fontsize=\scriptsize
}

% ---- Tiện ích khác ----
\usepackage{enumitem}     % tùy chỉnh danh sách
\usepackage{microtype}    % căn lề đẹp hơn (XeLaTeX hỗ trợ một phần)

% ---- Ngắt dòng và ngắt trang ----
% Tiếng Việt gần như không ngắt được từ nên TeX ít lựa chọn khi bẻ dòng; nới
% dung sai và cho phép giãn thêm ở lượt xếp cuối để không còn dòng tràn ra lề.
\tolerance=2000
\emergencystretch=3em
% Cấm để lại một dòng lẻ của đoạn ở đầu hoặc cuối trang.
\widowpenalty=10000
\clubpenalty=10000
\usepackage{float}        % [H] cho hình/bảng
\usepackage{subcaption}   % hình phụ (a), (b)...

% ---- Đánh số trang ----
% Số trang sẽ được bật trong main.tex sau phần mở đầu không đánh số
\pagestyle{plain}
